\section{Adjunction and special termination}\label{sec:special}

We use the generalized pairs and elementary steps of
Section~\ref{sec:relative}.  In particular, an elementary step is
projective over its own normal compact K\"ahler contraction base, is
negative for the adjoint on the negative side, and has a relatively
ample transformed adjoint on the positive side.  The models on which
elementary programs are run are globally strongly $\Q$-factorial.
The bases in a sequence need not be the same.  All boundaries and
higher line-bundle data in this section are rational.  The higher
data are nef and are carried by projective modifications.

Our goal is special termination: eventually both exceptional loci
avoid every generalized log canonical center.  Restriction to a center
reduces its proof to termination in a smaller dimension.  The center
need not satisfy the global strong factoriality condition, so we also
construct a crepant gdlt modification before running the smaller
program.

These assertions are proved in a fixed order.  Special termination
in dimension $d$ uses gdlt modifications and gdlt termination only in
dimensions less than $d$.  It then gives gdlt modifications in
dimension $d$.  For $d\leq3$, removing the floor and applying
Theorem~\ref{thm:gklt-three} gives gdlt termination in dimension $d$.
The induction ends with special termination and crepant gdlt
modifications in dimension four; it uses full gdlt termination only
through dimension three.
We first establish adjunction and discrepancy comparison on the
normalized centers.

\subsection{Adjunction on normalized strata}

We write $\operatorname{Nklt}(V,A+\bM)$ for the union of the
generalized log canonical centers of a glc pair.  A stratum below
means the closure of a coefficient-one stratum on the good snc open
set in the definition of gdlt.  Its normalization is always taken
before adjunction is iterated.

\begin{lemma}[Adjunction with a fixed higher denominator]
\label{lem:stratum-adjunction}
Let $(V,A+\bM)$ be a rational gdlt pair on a globally strongly
$\Q$-factorial normal compact K\"ahler space.  Suppose that $pM_W$
is a holomorphic line bundle on a projective carrier $W\to V$, for
an integer $p>0$.  If $T$ is the normalization of a generalized log
canonical center, then iterated divisorial adjunction along a
coefficient-one chain gives a rational gdlt pair
$(T,A_T+\bM^T)$ with effective boundary and an identity of rational
holomorphic adjoint bundles
\begin{equation}\label{eq:stratum-adjunction}
 K_T+A_T+M_T^T
   =\nu^*(K_V+A+M_V),
\end{equation}
where $\nu:T\to V$ is the natural map.  The b-line $\bM^T$ is
represented by restricting $\bM$ to the strict stratum on a higher
model; its denominator divides $p$.  Its generalized log canonical
centers map to generalized log canonical centers of $V$ properly
contained in $\nu(T)$.

More precisely, for a nonnegative coefficient set $I\subset[0,1]$
define
\begin{equation}\label{eq:adjunction-coefficients}
 \mathcal D_p(I)=
 \left\{
  1-\frac1r+\frac{\sum_{j=1}^{\ell}k_jd_j+k/p}{r}
  \ \middle|\ 
  \begin{array}{l}
  r\in\N_{>0},\ \ell,k,k_j\in\Z_{\geq0},\\
  d_j\in I,\ \sum_{j=1}^{\ell}k_jd_j+k/p\leq1
  \end{array}
 \right\}.
\end{equation}
If the coefficients of $A$ belong to $I$, those after $c$
successive adjunctions belong to $\mathcal D_p^{\circ c}(I)$,
where the superscript denotes the $c$-fold iterate.
If $I$ satisfies the descending chain condition (DCC), so does every
one of these sets.  No global
strong $\Q$-factoriality, or uniform Cartier index for $A_T$, is
asserted for $T$.
\end{lemma}

\begin{proof}
Choose a projective log resolution $h:W\to V$ carrying the nef
data which is an isomorphism at the general points of the good
strata under consideration.  We may replace $W$ further whenever
necessary, preserving these general points.  Since $V$ is globally
strongly $\Q$-factorial, the trace $M_V$ is a rational line bundle.
Lemma~\ref{lem:trace} and Lemma~\ref{lem:negativity} give
\begin{equation}\label{eq:adjunction-nef-defect}
 H=h^*M_V-M_W\geq0,
 \qquad H\text{ is }h\text{-exceptional}.
\end{equation}
Use compatible local canonical representatives to write
\[
 K_W+A_W+M_W=h^*(K_V+A+M_V),\qquad
 K_W+A_W^{\rm o}=h^*(K_V+A).
\]
These are actual rational sheaf comparisons, and
$A_W=A_W^{\rm o}+H$.

We first justify the ordinary companion in this construction.
The inequality above says that ordinary log discrepancies of
$(V,A)$ are at least the generalized log discrepancies.  In
particular $(V,A)$ is lc.  An ordinary zero-discrepancy place is
also a generalized zero-discrepancy place.  Its center therefore
meets the good open set.  Conversely, every generalized center
meets that set, where $H=0$, the data descend, and the log formula
is the ordinary snc formula.  Such a center is an ordinary
coefficient-one stratum.  Thus $(V,A)$ is ordinary dlt and its lc
centers are exactly the generalized ones.  This argument also
explains why forgetting the nef part does not introduce an
additional zero center outside the good open set; compare the
algebraic formulation in \cite[Remark~2.6(1)]{CT2023}.

Let $S\subset V$ be a prime component of $\lfloor A\rfloor$ and
let $S_W$ be its strict transform.  It is smooth after the chosen
resolution.  The coefficient of $S_W$ is one in both formulas.
The residue isomorphism restricts them to $S_W$ as
\begin{align*}
 (K_W+A_W+M_W)|_{S_W}
  &=K_{S_W}+(A_W-S_W)|_{S_W}+M_W|_{S_W},\\
 (K_W+A_W^{\rm o})|_{S_W}
  &=K_{S_W}+(A_W^{\rm o}-S_W)|_{S_W}.
\end{align*}
The normal space $S_W$ maps to the normalization $S^\nu$.
Push forward the indicated boundaries and the nef trace to
$S^\nu$.  Over its smooth large open set, the residue
identification is the adjunction identification of the canonical
sheaf.  It therefore defines, by reflexive extension, the
rational adjoint bundle on $S^\nu$ as the pullback of the
ambient adjoint.  This construction gives an actual line-bundle
isomorphism after a common positive multiple.  The prescribed
residue identification, and not equality of first Chern classes,
specifies the comparison.

The restricted log boundary on $S_W$ is snc with coefficients at
most one.  Thus the resulting generalized pair is sub-lc before
effectivity is checked.  Every zero stratum of that restricted
log formula is an intersection of $S_W$ with coefficient-one
components of the ambient log formula.  Its image is an ambient
generalized center and hence meets the good open set.  Over that
set the restriction is snc and its nef data descend.  The same
reasoning applies to the ordinary companion.  In particular,
after effectivity is established, both restricted pairs have the
appropriate dlt good-open property.  Their zero centers agree:
on a common carrier their log boundaries differ by the effective
restriction of $H$, while every generalized zero center meets
the open set where that restriction vanishes.

Here is the local effectivity and coefficient calculation, with
the issue of nonfactorial intermediate strata made explicit.
Suppose at some stage of the chain that $U$ is the normal space
already obtained, $C$ is its effective generalized boundary, and
$S\subset\lfloor C\rfloor$ is the next prime.  Carry along the
ordinary dlt companion just constructed.  Its floor has the same
strata.  Both adjoints are rational Cartier, although the
individual trace $M_U$ need not be rational Cartier on $U$.
Fix a prime divisor on $S^\nu$ whose coefficient is to be
computed, and work at an analytically general point of its image,
a codimension-two locus of $U$.

Take a general transverse surface slice there.  This operation
can be performed by fixing general values of local parameters
along that locus on a simultaneous embedded resolution.  We
choose the parameters transverse to the finitely many resolution
strata dominating the locus.  At a general point in question
the normal ambient space is Cohen--Macaulay: a normal space is
$S_2$, and its non-Cohen--Macaulay locus has codimension at least
three.  The sliced surface is $S_2$ and regular in codimension
one, hence normal.  The transverse restrictions of the
resolution formulas are crepant surface formulas, by
complete-intersection adjunction.  One can check this first
away from the distinguished point and then extend the
isomorphism of the rational Cartier adjoints.  In particular,
the ordinary surface pair is dlt along the coefficient-one
branch.  The boundary multiplicities computed on the slice are
the multiplicities at the original general point.

We use the following local surface fact.  An ordinary dlt
surface germ along a coefficient-one branch is either snc for
the reduced boundary, or is a plt germ with a smooth pair
chart modulo a small cyclic group of some order $r$.  The
branch in the chart is a coordinate curve.  The different of
that branch alone is $1-1/r$, and the local Cartier index of
any integral divisorial sheaf divides $r$.  This is the log
surface classification of
\cite[Theorem~4.15(1),(3), pp.~119--120, and Proposition~4.18]{KM1998}.
Its use here
is in the analytic-germ category: the numerical plt resolution
criterion gives a Hirzebruch--Jung string with the strict
boundary meeting an end; the corresponding normal analytic
surface germ is the cyclic quotient.  It is a classification
of the surface germ, and does not assert local factoriality
of the higher-dimensional space $U$.  In the snc case set
$r=1$.

Write the additional effective boundary on the slice as a sum
with coefficients $d_j$.  Pulling to the cyclic chart and
restricting to its smooth coordinate branch gives nonnegative
integral intersection multiplicities $k_j$.  Its contribution
to the different is $\sum_j k_jd_j/r$.  For the generalized
part, restrict a common higher model and the line bundle
representing $p\bM$ to the higher surface.  Define its trace on
the sliced surface by this proper pushforward; the restricted
higher log formula fixes the identification on the punctured
surface.  This avoids interchanging an arbitrary reflexive
hull with specialization.  The cyclic quotient fact makes
$rp$ times this trace Cartier.  Consequently its pullback
minus the higher nef bundle is an exceptional divisor with
denominator dividing $rp$.  It is effective by
Lemma~\ref{lem:negativity}, because the higher bundle is nef
on every curve of the projective surface resolution.  It
does not contain the strict coefficient-one branch, and its
restriction to that branch has coefficient $k/(rp)$ for
some $k\in\Z_{\geq0}$.  The resulting coefficient is therefore
\begin{equation}\label{eq:surface-generalized-different}
  1-\frac1r+\frac{\sum_jk_jd_j+k/p}{r}.
\end{equation}
This proves effectivity of the generalized different.
Omitting the nef correction proves effectivity of the ordinary
different carried along in the construction.  The coefficient
is at most one by the sub-lc log formula already established.
The local analytic adjunction and nef-defect calculation also
appears explicitly in the proof of
\cite[Lemma~3.1, pp.~5--6]{HXACC2023}; the local gdlt
modification in that proof is not used here as a global
modification theorem.

We can now repeat the construction.  On a resolution the chosen
chain is a chain of snc strata, and its successive restrictions
have the two log formulas just considered.  A zero stratum of
any restricted formula is a zero stratum of the preceding one;
its image is therefore an ambient center meeting the original
good open set.  This proves the center and gdlt assertions at
every stage.  Closures and normalizations introduce no finite
cover of the chosen center: each selected component is
generically isomorphic over it, so its normalization is its
unique normalization.  The restricted higher model is
projective over that normalization and is compact K\"ahler.
The higher line bundle remains nef, and its denominator still
divides $p$, under restriction and subsequent pullback.
This proves \eqref{eq:stratum-adjunction} and the asserted
description of the b-data.  A fixed chain on a common higher
model also shows that these b-data are compatible under any
birational comparison which preserves the general points of
the chain.

It remains to prove the DCC assertion rather than infer a
bound on all Cartier indices.  If $I$ is DCC, its positive
elements, when present, have a positive minimum.  The same is
true after adjoining $1/p$.  A sum of such elements bounded
by one consequently has a bounded number of nonzero summands.
Finite sums of a nonnegative DCC set satisfy DCC: from any
sequence of tuples of a fixed length one can successively
choose a subsequence on which each coordinate is
nondecreasing.  Thus the numerator sums in
\eqref{eq:adjunction-coefficients} satisfy DCC.  For bounded
$r$, their images in that formula satisfy DCC.  In a strictly
decreasing sequence of displayed coefficients with first
term $c_0<1$, the inequality
\[
  1-\frac1r\leq c\leq c_0
  \quad\Longrightarrow\quad r\leq\frac1{1-c_0}
\]
bounds $r$.  A sequence starting at one reduces to this case
after its first strict decrease.  Therefore
$\mathcal D_p(I)$ is DCC, and induction proves the same for
each finite iterate.
\end{proof}

\subsection{Comparison on a stratum}

Adjunction supplies an effective generalized boundary with coefficients
in a fixed DCC set.  To use those coefficients along a sequence, we need
monotonicity of the restricted discrepancies.  Strictness must also
detect an ambient exceptional intersection of arbitrary codimension
on the stratum.

\begin{lemma}[Restricted discrepancy comparison]
\label{lem:stratum-comparison}
Suppose that
\[
 (V,A+\bM)\dashrightarrow (V^+,A^++\bM)
\]
is a small elementary gdlt step over $Z$.  Let $T$ be a
normalized generalized log canonical center whose general
point and coefficient-one chain are unchanged by the step,
and let $T^+$ denote the corresponding normalization on the
positive side.  Put
$D_T=K_T+A_T+M_T^T$ and $D_{T^+}=K_{T^+}+A_{T^+}+M_{T^+}^T$.
Both strata map projectively and bimeromorphically to the
normalization $S$ of their common image in $Z$.  The adjoints
are respectively antiample and ample over $S$.

On a common smooth higher model, with projective maps $p$ and
$q$ to $T$ and $T^+$, respectively,
\begin{equation}\label{eq:restricted-positive-difference}
 p^*D_T-q^*D_{T^+}=F_T\geq0.
\end{equation}
For every prime divisor $E$ on a proper bimeromorphic analytic model
of the stratum, the generalized discrepancies do not decrease.
They increase strictly if its center on either stratum is
contained in the corresponding ambient exceptional locus.
\end{lemma}

\begin{proof}
The restricted morphisms factor through $S$ by normality.
They are projective, since restriction and this finite
factorization preserve relative ampleness; they are
bimeromorphic because the ambient step is an isomorphism
at the general point of the stratum.  Restriction of the
relatively ample adjoint, and then pullback by the
normalization, gives the asserted signs.

Take a common projective log resolution of the ambient step,
carrying its common nef data and preserving the general
points of the chosen snc chain.  If $a,b$ are its projections,
Lemma~\ref{lem:negativity} gives the actual divisor comparison
\[
 F=a^*(K_V+A+M_V)-b^*(K_{V^+}+A^++M_{V^+})\geq0.
\]
This difference is anti-nef over $Z$.  Its support contains
the full fibers over the non-isomorphism locus in $Z$:
opposite relative ampleness makes the difference nontrivial
over each such point, and the fiber-support assertion of
Lemma~\ref{lem:negativity} applies.  The strict transform of
the chosen stratum is not contained in $\Supp F$, since its
general point lies in the common isomorphism open set.

Restrict the two log formulas along the common snc chain.
At each restriction the residue canonical bundles are the
same, and the common higher nef bundles are the same.
They cancel.  Each restricted difference is effective,
because the next strict stratum is not contained in its
support.  At the last stage this gives
\eqref{eq:restricted-positive-difference} as an equality
with its prescribed meromorphic identification.  The
coefficient of $E$ in this difference is the increase of its
generalized log discrepancy.

For strictness, first make $E$ divisorial on a smooth higher
stratum.  This can be achieved by projective blowups of its
successive centers; perform the same blowups in the ambient
resolution and resolve further, preserving the generic
chain.  If the center of $E$ is in the ambient exceptional
locus, its entire inverse image in the common ambient model
lies in $\Supp F$, by the full-fiber assertion.  Therefore
the center on the higher stratum lies in the support of the
restricted effective Cartier divisor.  Its defining local
equation has positive order at $E$.  This proves strictness
on either side.  In particular, the argument detects an
intersection in arbitrary codimension on the original
stratum; it does not require that intersection already to
be a divisor there.  The algebraic counterparts are
\cite[Lemma~2.8]{LMT2022} and
\cite[Proposition~3.2]{Moraga2020}; the argument above uses
the analytic projective negativity statement directly.
\end{proof}

\subsection{The triangular induction}

We state the three assertions together to make their dependency order
explicit.  The first stops steps near the generalized log canonical
centers; the second constructs crepant models; the third supplies the
lower-dimensional termination used in the next induction stage.

\begin{theorem}[Special termination]\label{thm:special}
Let $d\leq4$, and let
\[
 (V_0,A_0+\bM)\dashrightarrow(V_1,A_1+\bM)
 \dashrightarrow\cdots
\]
be a sequence of elementary steps for a rational gdlt pair
on globally strongly $\Q$-factorial normal compact K\"ahler
spaces of dimension $d$, with fixed higher nef b-line data.
There exists $i_0$ such that, for every $i\geq i_0$, the
exceptional loci on both sides of the step are disjoint
from all generalized log canonical centers on those sides.
No common contraction base for the sequence is assumed.
\end{theorem}

\begin{proposition}[Crepant gdlt modification]
\label{prop:gdlt-mod}
Let $(V,A+\bM)$ be a rational glc pair of dimension at most
four on a normal compact K\"ahler space.  Its higher nef
line data are carried by a projective modification, and
$K_V+A+M_V$ is a rational holomorphic line bundle.
There is a projective bimeromorphic morphism $h:Y\to V$
such that $Y$ is normal compact K\"ahler and globally
strongly $\Q$-factorial, $(Y,A_Y+\bM)$ is gdlt with
$A_Y\geq0$, and
\[
 K_Y+A_Y+M_Y=h^*(K_V+A+M_V).
\]
Every $h$-exceptional prime has coefficient one in $A_Y$,
or equivalently has generalized log discrepancy zero over
$(V,A+\bM)$.  Global strong $\Q$-factoriality is not assumed
for $V$.
\end{proposition}

\begin{theorem}[Gdlt termination in dimensions at most three]
\label{thm:gdlt-three}
Every sequence as in Theorem~\ref{thm:special} in dimension
at most three is finite.  The assertion concerns arbitrary
choices of the permitted elementary steps, and does not
assume pseudo-effectivity of the adjoint.
\end{theorem}

\begin{proof}[Joint proof of Theorem~\ref{thm:special},
Proposition~\ref{prop:gdlt-mod}, and Theorem~\ref{thm:gdlt-three}]
For each $d$ in increasing order we prove special
termination, then the modification assertion, and, if
$d\leq3$, gdlt termination.  In dimension zero these
statements are immediate.  Fix $d>0$, assume all the
appropriate assertions in smaller dimensions, and first
prove special termination in dimension $d$.

\smallskip
\noindent\emph{Stabilizing centers.}
There are only finitely many divisorial steps, by
Lemma~\ref{lem:cycle-loss}; hence we may discard a finite
prefix and consider only small steps.  Generalized
discrepancies do not decrease.  A zero-discrepancy place on
the positive side was consequently a zero-discrepancy
place before the step, and strictness in
Lemma~\ref{lem:negativity} excludes its center from either
exceptional locus.  Thus every new center is the transform
of an old center with unchanged general point.  If the
exceptional locus contains the general point of an old
center, that center disappears.  There are finitely many
centers on the initial model, so after another finite
prefix their list is unchanged and all steps are
isomorphisms at their general points.

For each center choose a chain of coefficient-one strata
at its general point.  Transform this choice along the
sequence.  Lemma~\ref{lem:stratum-adjunction} gives the
corresponding normalized strata $T_i$ and their adjoints
$D_{T_i}$.  The b-line data on these strata are the same
under their birational identifications, with one fixed
higher denominator.  The ambient boundary coefficients
belong to a fixed finite set.  The coefficients after
adjunction therefore belong to one fixed DCC set, depending
only on that set, the denominator, and the length of the
chosen chain.

We now induct on the dimension $e$ of a surviving center
to show eventual disjointness.  A point center is already
disjoint after the stabilization just made.  For a center
of positive dimension $e$, its proper zero substrata have
smaller dimension and there are finitely many of them.
Discarding a further finite prefix, we may assume that
both exceptional loci avoid every such substratum at
every remaining step.  In particular, the induced maps
on $T_i$ are isomorphisms in neighborhoods of their
non-gklt loci.

\smallskip
\noindent\emph{Finitely many extractions on the stratum.}
Identify divisorial places over the $T_i$ through common proper
bimeromorphic analytic models.  If a prime
divisor $E$ is extracted by $T_i\dashrightarrow T_{i+1}$,
then $E$ is a divisor on $T_{i+1}$ and its center there is
contained in the ambient exceptional locus.  Otherwise
the ambient isomorphism near its general point would
provide its strict transform as a divisor on $T_i$.
Lemma~\ref{lem:stratum-comparison} gives
\[
 a(E;T_i,A_{T_i}+\bM^T)
 <a(E;T_{i+1},A_{T_{i+1}}+\bM^T)\leq1,
\]
where the last inequality uses effectivity of the boundary.
By monotonicity its discrepancy on $T_0$ is less than one.
Its center on $T_0$ cannot be contained in
$\operatorname{Nklt}(T_0,A_{T_0}+\bM^T)$: the maps are
isomorphisms near that locus and its transforms, so a
valuation centered there cannot acquire a divisorial
center through the transformations under consideration.

There are only finitely many such valuations by the
off-non-gklt assertion of Lemma~\ref{lem:low-places}.
In addition to its exceptional valuations, include the
finitely many divisors already on $T_0$ having discrepancy
less than one; these are components of its positive
boundary.  For any one of this finite set, the coefficients
on the models where it is extracted form a strictly
decreasing sequence: discrepancy never decreases, and
each new extraction gives a strict increase.  Those
coefficients are in the fixed DCC set.  Thus each
valuation is extracted only finitely often.  Altogether
there are finitely many extractions.

\smallskip
\noindent\emph{Removing divisorial changes and stabilizing coefficients.}
After these extractions, the induced bimeromorphic maps
extract no divisors.  Lemma~\ref{lem:cycle-loss}, applied
to the span of prime $(e-1)$-cycles of the compact
K\"ahler strata, shows that only finitely many further
maps contract a divisor.  In dimension one a proper
bimeromorphic map between normal curves is already an
isomorphism.  We may therefore assume that every induced
map is an isomorphism in codimension one.

With primes matched on this tail, their boundary
coefficients do not increase, by
Lemma~\ref{lem:stratum-comparison}.  No component with
coefficient zero can acquire a positive coefficient.
Only the finite boundary support at the start of the
tail needs to be considered, and DCC stabilizes each
coefficient.  Consequently
\[
 (T_i\dashrightarrow T_{i+1})_*A_{T_i}=A_{T_{i+1}}.
\]
Let $S_i$ be the normalization of their common image in
the contraction base.  Neither $T_i\to S_i$ nor
$T_{i+1}\to S_i$ contracts a prime divisor.  Indeed such
a prime is matched in codimension one on the other
stratum, but its center is in the corresponding ambient
exceptional locus.  Strict discrepancy increase would
contradict coefficient stabilization.

We have thus obtained small diagrams over the varying
$S_i$, with antiample and ample adjoints, respectively,
and strict-transform boundary and b-data.  The effective
difference of Lemma~\ref{lem:stratum-comparison} is now
exceptional over both stratum models: its coefficient
at a divisor on either model is zero by the matched data.
If one of the induced birational maps is an isomorphism,
the matched boundary and b-line identify its adjoints
as actual rational bundles with their canonical
meromorphic comparison.  The restricted difference is
then zero on a common model.  The strictness assertion
shows that this step has no ambient exceptional
intersection with the stratum.  Thus only nontrivial
small stratum surgeries remain to be excluded.

\smallskip
\noindent\emph{Lifting the small surgeries in dimension $e<d$.}
The induction hypothesis in dimension $e$ gives a
projective crepant gdlt modification
$h_0:Y_0\to T_0$ with globally strongly
$\Q$-factorial source.  Over $S_0$ run the relative
elementary steps of Proposition~\ref{prop:relative-steps}
for its crepant adjoint.  By gdlt termination in
dimension $e$, this run is finite and ends at a model
$Y_1$ whose adjoint is nef over $S_0$.

On a common smooth model of the initial lift, its
endpoint, and $T_1$, write the actual comparisons as
\[
 P_0=P_1+G=P_++F.
\]
Here $P_0$ pulls back the initial crepant adjoint,
$P_1$ pulls back the adjoint of $Y_1$, and $P_+$ pulls
back $D_{T_1}$.  The relative steps give $G\geq0$
exceptional over $Y_1$.  The restricted comparison
just established, pulled through the initial crepant
modification, gives $F\geq0$ exceptional over $T_1$.
Lemma~\ref{lem:nef-end-comparison} applies: $P_1$ is
nef over $S_0$ and $D_{T_1}$ is ample there.  It gives
$P_1=P_+$ and a projective crepant morphism
$h_1:Y_1\to T_1$.

This endpoint is again a gdlt modification of the
required type.  The elementary steps preserve gdlt
and effectivity.  They extract no prime divisor, and
the stratum map below is small.  Any prime on $Y_1$
exceptional over $T_1$ was therefore exceptional on
$Y_0$ over $T_0$.  If it survives, its coefficient is
still one, since boundaries are pushed forward.
Crepancy identifies that coefficient with its
zero-discrepancy extraction over $T_1$.

Repeat over $S_1,S_2,\ldots$.  Each nontrivial stratum
surgery forces at least one lifted elementary step.
Indeed its negative morphism has a contracted curve
$C$ with $D_{T_i}\cdot C<0$.  Projectivity gives a
curve on the lift dominating $C$ with positive
degree, on which the pulled-back adjoint is negative.
The lifted model is therefore not already nef over
$S_i$.  Infinitely many nontrivial stratum surgeries
would give an infinite concatenation of elementary
gdlt steps in dimension $e$.  This contradicts the
lower-dimensional gdlt termination hypothesis, which
allows varying compact contraction bases.  There
are only finitely many such surgeries.  The preceding
isomorphism case now proves disjointness for the
chosen center.

This completes the induction on $e$.  Taking the
maximum of the finitely many resulting indices gives
special termination in dimension $d$, on both sides
of every remaining step.  This proves
Theorem~\ref{thm:special} in that dimension.

\smallskip
\noindent\emph{Constructing the modification in dimension $d$.}
Now let $(V,A+\bM)$ be an arbitrary rational glc pair
as in Proposition~\ref{prop:gdlt-mod} in dimension $d$.
Take a projective log resolution $h:W\to V$ with $W$ smooth,
carrying the higher nef bundle.  Let
\[
 \Gamma=h_*^{-1}A+\sum_{E\subset\Exc(h)}E.
\]
The pair $(W,\Gamma+\bM)$ is gdlt: its boundary is
snc and its data descend to $W$.  Its adjoint has
the actual comparison
\begin{equation}\label{eq:gdlt-modification-excess}
 K_W+\Gamma+M_W=h^*(K_V+A+M_V)+Q,
 \qquad
 Q=\sum_{E\subset\Exc(h)}
       a(E;V,A+\bM)E\geq0.
\end{equation}
The support of $Q$ is contained in $\lfloor\Gamma\rfloor$.
Run relative elementary steps over $V$ using
Proposition~\ref{prop:relative-steps}.  On every model
the pushforward of \eqref{eq:gdlt-modification-excess}
remains the identity of its adjoint with the pulled-back
original adjoint plus the effective exceptional
transform of $Q$.  Every surviving component of that
transform is a floor component.

If this run were infinite, special termination in
dimension $d$, just proved, would make its tail
disjoint from the floor.  A contracted curve on such
a tail lies outside the support of $Q$ and has
nonnegative degree against this effective rational
Cartier divisor.  Its degree against the pulled-back
original adjoint is zero.  This contradicts negativity
of the elementary step.  Thus the run is finite;
continuation ensures it ends at relative nefness.
On the endpoint $Y$, the remaining transform $Q_Y$
is nef over $V$ and exceptional, so
Lemma~\ref{lem:negativity} gives $Q_Y\leq0$.  It is
also effective, hence zero.  The comparison is
crepant.  The resulting model is projective over $V$,
compact K\"ahler, globally strongly $\Q$-factorial,
and gdlt, by the relative construction.  Any prime
exceptional over $V$ is the transform of an
exceptional prime on $W$, and has coefficient one
in the pushed boundary.  This proves
Proposition~\ref{prop:gdlt-mod} in dimension $d$.

\smallskip
\noindent\emph{Gdlt termination in dimension $d\leq3$.}
Finally suppose $d\leq3$ and consider any infinite
elementary gdlt sequence in this dimension.  Discard
its finite divisorial prefix and apply special
termination.  On the resulting tail all exceptional
loci avoid $\lfloor A_i\rfloor$.  Replace its boundary
by $A_i-\lfloor A_i\rfloor$ on each model, retaining
the same b-line.  The resulting pairs are effective
and gklt.  To check the latter, discrepancies cannot
decrease on removing an effective rational Cartier
floor.  Each old zero-discrepancy place has center
in the floor, and the order of its pullback is
positive, so its discrepancy becomes positive.
Equivalently this is the snc calculation on a fixed
carrier and the gdlt good-open condition.

The floor has degree zero on every contracted curve
on both sides, because it is disjoint from the
exceptional loci.  Thus all steps of this tail are
still negative and positive, respectively, for the
new adjoints.  Their boundaries are strict transforms,
their nef data are unchanged, and their one-ray
condition for global line-bundle degrees is
unchanged.  They are consequently elementary gklt
steps in the category of
Theorem~\ref{thm:gklt-three}.  That theorem excludes
an infinite tail.  This proves
Theorem~\ref{thm:gdlt-three} in dimension $d$ and
completes the triangular induction through
dimension four.
\end{proof}

\begin{remark}
The proof constructs auxiliary strong-$\Q$-factorial
models only after adjunction, or over the space to
which a gdlt modification is requested.  It never
assumes that a normalized stratum is locally
$\Q$-factorial.  Its uses of curve-nefness occur over
individual projective bimeromorphic morphisms;
the termination assertion concerns the entire
sequence of compact models with varying bases.
No analytic special-termination theorem is deduced
solely by citing its projective counterpart.
\end{remark}
